% Compile with pdfLaTeX.
% The helvet package supplies a Helvetica-compatible sans-serif family.
\documentclass[10pt,a4paper]{article}

\usepackage[a4paper,top=10.5mm,bottom=10.5mm,left=12mm,right=12mm]{geometry}
\usepackage[T1]{fontenc}
\usepackage{helvet}
\usepackage{fontawesome5}
\renewcommand{\familydefault}{\sfdefault}
\usepackage{xcolor}
\usepackage{tabularx}
\usepackage{enumitem}
\usepackage{hyperref}
\usepackage{microtype}
\usepackage{ragged2e}

\definecolor{cvblue}{HTML}{164F8C}
\hypersetup{
  colorlinks=true,
  urlcolor=black,
  linkcolor=black,
  pdfauthor={Zander Polson},
  pdftitle={Zander Polson CV}
}
\hyphenpenalty=10000
\exhyphenpenalty=10000

\pagestyle{empty}
\setlength{\parindent}{0pt}
\setlength{\parskip}{0pt}
\setlength{\tabcolsep}{0pt}
\renewcommand{\arraystretch}{0.95}
\raggedbottom
\emergencystretch=1em

\newcommand{\cvsection}[1]{%
  \vspace{0.90em}%
  {\color{cvblue}\bfseries\fontsize{11.2}{11.8}\selectfont #1}\par
    \vspace{-0.5em}%
  {\color{black!58}\rule{\linewidth}{0.35pt}}\par
  \vspace{0.38em}%
}

\newcommand{\job}[4]{%
  \begin{tabularx}{\linewidth}{@{}>{\raggedright\arraybackslash}X>{\raggedleft\arraybackslash}p{37mm}@{}}
    \bfseries #1 & \bfseries #2 \\
    \itshape #3 & \itshape #4
  \end{tabularx}\par
  \vspace{-0.10em}%
}

\newcommand{\project}[1]{%
  \vspace{0.18em}%
  {\bfseries #1}\par
  \vspace{-0.12em}%
}

\newcommand{\sidebarheading}[1]{%
  \vspace{0.50em}%
  {\bfseries #1}\par
  \vspace{0.08em}%
}

\newcommand{\education}[4]{%
  {\bfseries\fontsize{9.9}{10.5}\selectfont #1}\par
  {\bfseries\fontsize{9.9}{10.5}\selectfont #2}\par
  {\itshape #3}\par
  {\itshape #4}\par
  \vspace{0.52em}%
}

\setlist[itemize]{
  leftmargin=1.05em,
  label=\textbullet,
  labelsep=0.30em,
  topsep=0pt,
  partopsep=0pt,
  parsep=0pt,
  itemsep=0.05em
}

\begin{document}
\fontsize{10}{10.8}\selectfont

% Header
{\color{cvblue}\fontsize{20.5}{21}\selectfont ZANDER POLSON}\par
\vspace{0.12em}
{\fontsize{7.9}{8.6}\selectfont Computer Engineer | Robotics, Embedded Systems and Computer Vision}\par
\vspace{0.05em}
{\fontsize{7.8}{8.5}\selectfont
	\faWhatsapp\ +27 749248224 |
	\href{mailto:zander@polsons.info}{\underline{zander@polsons.info}} |
	\href{https://github.com/GitM3}{\underline{github.com/GitM3}}}\par

\vspace{0.05em} %

% Main two-column structure. Content is laid out with normal LaTeX flow;
% there is no absolute positioning or coordinate-based sentence placement.
\begin{minipage}[t]{0.686\textwidth}
	\RaggedRight
	\cvsection{PROFILE}
	Computer engineer and Advanced Robotics master’s student interested in embedded
	systems, robotics, computer vision, and system integration.
	Experience in startup R\&D, helping develop embedded IoT and perception systems
	from prototype to client-facing demonstrations.
	Outside of engineering, I pursue fitness and various adventure-sports.

	\cvsection{EXPERIENCE}
	\job{Chiba Institute of Technology}{Sep 2025 - Present}{Master's Program in Advanced Robotics}{Japan}
	\begin{itemize}
		\item Currently pursuing graduate studies in Robotics, with research focusing on synthetic data creation for agricultural robotic applications.
		\item Presented work at RoboMech2026 conference on scalable sweet-pepper plant generation pipeline using Blender and generative AI models.
	\end{itemize}

	\vspace{0.34em}%%
	\job{UViRCO}{Nov 2024 - Jun 2025}{Junior Embedded Systems Engineer}{South Africa / China}
	\begin{itemize}
		\item Helped take the platform from prototype to first deployment in less than six months, including an on-site demonstration for major clients in China.
		\item Developed an embedded IoT and computer-vision platform for Raspberry Pi Compute Module 4/5 using C/C++, Python, and OpenCV, delivering a functional prototype capable of real-time image analysis.
		\item Integrated Redis-based microservices, GStreamer video pipelines, and an ONVIF C server into a Docker Compose deployment, enabling scalable video streaming and standardized device communication.
		\item Built CI/CD workflows with GitHub Actions and contributed a Python remote-update service for edge devices, reducing deployment time and ensuring reliable over-the-air updates.
	\end{itemize}
	\vspace{0.34em}%%

	\job{Rapid Mobile}{Jun 2024 - Jul 2024}{Software Engineering Intern}{South Africa}
	\begin{itemize}
		\item Implemented a web-interface feature in a custom Qt-based desktop application using C++ and QWebEngine, enabling client to access a remote radio interface in real-time.
	\end{itemize}
	\vspace{0.34em}%%

	\job{Various organisations}{Feb 2020 - Nov 2020}{Community and Hospitality Volunteer}{Japan}
	\begin{itemize}
		\item Volunteered at Tannenhof ski-resort hotel in Nozawa Onsen, where I assisted the owner with general maintenance, guest check-ins and international guest tours to ensure a welcoming and memorable experience.
		\item Lived and worked with the Dana Village community and Chardjou-sol organic tomato farm in Fukushima, supporting daily farm operations such as planting and harvesting while learning about no-irrigation farming techniques.
		\item Volunteered at Hisaeda Udon restaurant in Takamatsu, gaining practical experience in a Japanese-speaking work environment.
	\end{itemize}

	\cvsection{PROJECTS}
	\project{B.Eng. Final-Year Project -- Real-Time Mobile SLAM }
	\begin{itemize}
		\item Developed a real-time Android SLAM application using Qt and OpenCV, with visual odometry as a fallback localization mode.
	\end{itemize}

	\project{Gauteng Department of Roads and Transport -- Emission station IoT Prototype }
	\begin{itemize}
		\item Developed and delivered a PCB-based, solar-powered IoT prototype for monitoring environmental pollution, using GSM connectivity for remote data transmission, through the University of Pretoria’s Joint Community-Based Project.
	\end{itemize}

	\project{Team Project -- RF-DETR Bottle-Following Robot}
	\begin{itemize}
		\item Developed the perception layer of a bottle-following mobile robot using RF-DETR on an NVIDIA Jetson Orin. Collaborated with other Japanese graduates who developed the LiDAR-based navigation stack. We integrated the perception and navigation components using ROS to enable real-time target detection and following.
	\end{itemize}

	\cvsection{AWARDS}
	\begin{itemize}
		\item \textbf{UP Robot Race Day:} Finalist in the University of Pretoria EECE Department's 18-bit microcontroller-based autonomous robotic vehicle race.
	\end{itemize}
\end{minipage}%
\hfill
\begin{minipage}[t]{0.279\textwidth}
	\RaggedRight
	\cvsection{EDUCATION}
	\education{MSc. Advanced Robotics}{2025 - Present}{Chiba Institute of Technology}{Japan}
	\education{B.Eng. Computer Engineering}{2021 - 2024}{University of Pretoria}{South Africa}

	\cvsection{SKILLS}
	\sidebarheading{Programming}
	C/C++\par
	Python\par
	Qt/QML\par
	MATLAB\par
	Git\par

	\sidebarheading{Embedded Devices}
	STM32\par
	Raspberry Pi CM4/CM5\par
	NVIDIA Jetson \par
	Intel Altera SoC FPGA\par
	PIC18\par

	\sidebarheading{Languages}
	English\par
	Japanese (JLPT N3)\par
\end{minipage}

\end{document}
